\section{Introduction}\label{sec:introduction}
For the past two and a half million years, Earth's temperature has fluctuated. Geological records show that global temperatures oscillate between warmer and colder temperatures with a period of about 100,000 years \parencite{herringHasntEarth2020}. However, this pattern cannot be described with a simple sinusoidal relationship; it shows slow, irregular cooling that appears almost jagged on a graph, followed by rapid deglaciation that forces ice back to the poles. What drives this rhythm is one of the most important questions in paleoclimatology.

The prevailing theory involves Milankovitch cycles. These periodic variations in Earth's orbital parameters (eccentricity, axial tilt, and precession) alter the distribution of insolation. Yet this model has its flaws; the temperature differences the Earth experiences are far too great to be explained by these cycles alone. This mismatch catalyzed the search for internal mechanisms that could explain the oscillation. These are processes within the climate system itself that could amplify external forces or even drive them autonomously.

Energy Balance Models (EBMs) offer a mathematical framework for this search. Rather than resolving the full complexity of the atmosphere and oceans, an EBM simplifies the problem to a few variables governed by the law of conservation of energy \parencite{lohmannTemperaturesEnergy2020}. The one-variable Budyko-Sellers model (1969) successfully reproduces multiple climate equilibria associated with glaciation and deglaciation \parencite{budykoEffectSolar1969, sellersGlobalClimatic1969}. However, a single-variable ODE cannot produce oscillations. \footnote{Explaining these cycles requires at least two variables.}

The K\"allen-Crafoord-Ghil model (1979), or KCG model, achieves this by relating temperature $T$ to meridional ice-sheet extent $L$. The resulting system can produce a wider variety of behaviours, including the possibility that a stable steady state can lose stability and give rise to self-sustained oscillations \parencite{kallenFreeOscillations1979}.

This essay will explore the following research question:\\
\textit{How far from equilibrium does the linearization of a two-variable energy-balance climate model remain a valid predictor of the system's behaviour?}

To answer this question, we will first build the model from the pre-existing literature, then linearize it, and finally test it numerically. The model being tested is not novel to this essay; linearization is a widely used technique for predicting the behaviour of a system of coupled differential equations. Instead, it will examine the deviation of the approximation from the actual system to better assess the reliability of these methods.

\subsection{Overview of Paper}
Section \ref{sec:background} defines the mathematics used throughout the paper and begins building the model using $\tanh$ albedo feedback, which is expanded to a cubic function and then rescaled before being coupled with the ice equation of the KCG model. Section \ref{sec:linearizing} linearizes this system about its equilibrium, and Section \ref{sec:auxiliary-equation} solves the resultant second-order ODE to predict the period, $P$, and growth rate, $\alpha$, of small oscillations. Section \ref{sec:numerical-analysis} solves the entire nonlinear system utilizing RK4 and measures how far the actual trajectory differs from the prediction as the control parameter, $a$, and the starting displacement, $\eta_0$, vary. Section \ref{sec:results-discussion} interprets the results, defines a radius of validity, answers the research question, and compares it with existing literature. Section \ref{sec:evaluation-conclusion} evaluates the method's limitations, proposes extensions, and concludes the paper.